\section{Introduction}
\label{sec:introduction}

Object detectors are increasingly deployed under simultaneous accuracy,
latency, memory, and energy constraints. Post-training quantization (PTQ)
reduces numerical precision without repeating full model training and is
therefore attractive for rapid deployment\citep{jacob2018integer,banner2019post}.
Integer arithmetic remains a standard path to compact inference, while FP8
provides an 8-bit floating-point alternative with a wider dynamic range and
increasing accelerator support\citep{micikevicius2022fp8,shen2024fp8}. These
benefits matter most outside the laboratory, where images are also affected by
sensor noise, motion, visibility loss, and lossy transmission. Common-corruption
benchmarks and detector-specific studies show that such changes can substantially
reduce recognition and localization accuracy\citep{hendrycks2019corruptions,
michaelis2019detectorrobustness,liu2024realworldcorruptions}.

Quantization and corruption are usually assessed with different summaries.
Quantization papers emphasize a clean full-precision-to-quantized accuracy gap,
whereas corruption papers emphasize an absolute clean-to-corrupted drop. A raw
corrupted-image difference between two quantized engines answers neither
question cleanly. If FP8 already exceeds INT8 on the matched clean input, the
corrupted FP8--INT8 gap contains that pre-existing difference plus any additional
format--corruption interaction. Calling the raw gap a robustness advantage
therefore confounds clean fidelity with corruption-specific behavior.

A second ambiguity arises near an accuracy floor. Severe corruptions may drive
both engines toward very low AP, causing their difference to contract. A signed
interaction can then appear favorable precisely because neither engine remains
useful. A defensible evaluation must therefore distinguish three claims:
(i) matched-clean fidelity, (ii) absolute corrupted accuracy, and (iii) the
change in the format discrepancy after the matched-clean difference has been
removed. These claims are related, but none substitutes for the others.

The labels ``INT8'' and ``FP8'' also require care. An executable TensorRT engine
is a treatment comprising number format, quantizer placement, calibration,
graph rewriting, mixed-precision coverage, runtime tactics, preprocessing, and
decoding. Two engines with the same nominal format can behave differently when
any of these components changes. Our target is consequently the comparison of
the recorded TensorRT INT8 and FP8 treatments, not an intrinsic ordering of all
possible INT8 and FP8 implementations.

We introduce a paired, floor-guarded protocol that makes these distinctions
explicit. For every clean or corrupted condition, images are materialized once
and then consumed byte-for-byte by both quantized engines. The direct estimand
is a four-cell difference in differences,
\begin{equation*}
\Delta E=
(\mathrm{FP8}-\mathrm{INT8})_{\mathrm{corrupt}}-
(\mathrm{FP8}-\mathrm{INT8})_{\mathrm{clean}},
\end{equation*}
computed within common image-bootstrap draws. An accompanying absolute-AP
guard prevents gap contraction near a shared failure floor from being labelled
robustness. The exploratory landscape contains four datasets, three YOLO11
capacity rungs, four synthetic corruption families, and three ordinal
severities, for 144 direct comparison cells. We then probe the result with
untouched VOC/KITTI final partitions, alternative weights, repeated corruption
materializations, training and calibration seeds, a fixed-class-universe
bootstrap, measured engine runtime, and portability stress cases on RT-DETR and
RetinaNet.

The main contributions are:
\begin{itemize}[leftmargin=1.45em,itemsep=2pt,topsep=3pt,parsep=0pt]
  \item a byte-matched four-cell estimand that isolates the corruption-induced
  change in the INT8--FP8 discrepancy while preserving image-level covariance;
  \item a floor-guarded reporting rule that jointly presents clean fidelity,
  absolute corrupted AP, and the clean-adjusted interaction instead of reducing
  them to one robustness score;
  \item a decision-impact analysis showing when the clean adjustment changes
  the practical interpretation of the raw corrupted gap; and
  \item a layered validation study that makes the conditional scope of split,
  realization, seed, class-universe, runtime, and architecture evidence explicit.
\end{itemize}

The empirical conclusion is deliberately narrower than a format ranking. FP8
preserves more matched-clean AP in every YOLO dataset--model block, but the
corruption interaction is heterogeneous and changes with the evaluation scope.
The contribution is therefore an evaluation methodology and an auditable body
of treatment-conditional evidence, not a universal claim that one 8-bit format
is more robust.
